\chapter{Вероятност и~вероятностно пространство}

\section{Дефиниция и~аксиоми}

\begin{defbox}[~(Вероятност)]
Нека $\Omega$ е пространство от елементарни събития и~нека
$\mathcal{A}$ е $\sigma$-алгебра над $\Omega$. \textbf{Вероятността}
(вероятностна мярка) е функция $\mathbb{P}:\mathcal{A}\to[0,1]$, за
която:
\begin{itemize}[nosep]
  \item $\mathbb{P}(A) \geq 0$ за всяко $A\in\mathcal A$;
  \item $\mathbb{P}(\Omega) = 1$;
  \item ако $A_1,\dots,A_n,\dots$ са несъвместими (т.е.\ $A_i\cap A_j=\varnothing$ за $i\neq j$), то
  $\displaystyle \mathbb{P}\Big(\bigcup_{i=1}^{\infty}A_i\Big) = \sum_{i=1}^{\infty}\mathbb{P}(A_i)$ ($\sigma$-адитивност).
\end{itemize}
\end{defbox}

\begin{corbox}[~(основни свойства на вероятността)]
Нека $\mathbb{P}:\mathcal{A}\to[0,1]$ е вероятност. Тогава са в~сила:
\begin{enumerate}[label=\alph*)]
  \item $0 \leq \mathbb{P}(A) \leq 1$;
  \item $\mathbb{P}(\varnothing) = 0$;
  \item $\mathbb{P}(\bar A) = 1-\mathbb{P}(A)$;
  \item ако $A\cap B=\varnothing$, то $\mathbb{P}(A\cup B)=\mathbb{P}(A)+\mathbb{P}(B)$;
  \item в~общия случай $\mathbb{P}(A\cup B) = \mathbb{P}(A)+\mathbb{P}(B)-\mathbb{P}(A\cap B)$;
  \item ако $A\subseteq B$, то $\mathbb{P}(A)\leq\mathbb{P}(B)$ (монотонност);
  \item (непрекъснатост отдолу) ако $A_1\subseteq A_2\subseteq\dots\subseteq A_n\subseteq\dots$, то
  $\displaystyle \mathbb{P}\Big(\bigcup_{j=1}^{\infty}A_j\Big)=\lim_{n\to\infty}\mathbb{P}(A_n)$;
  \item (непрекъснатост отгоре) ако $A_1\supseteq A_2\supseteq\dots\supseteq A_n\supseteq\dots$, то
  $\displaystyle \mathbb{P}\Big(\bigcap_{i=1}^{\infty}A_i\Big)=\lim_{n\to\infty}\mathbb{P}(A_n)$;
  \item (изброима субадитивност) ако $A_i\in\mathcal A$ за $i\geq 1$, то
  $\displaystyle \sum_{i=1}^{\infty}\mathbb{P}(A_i) \geq \mathbb{P}\Big(\bigcup_{i=1}^{\infty}A_i\Big)$,
  като същото важи и~за крайно обединение: $\sum_{i=1}^{n}\mathbb{P}(A_i)\geq\mathbb{P}\big(\bigcup_{i=1}^{n}A_i\big)$ за $n\in\mathbb N$.
\end{enumerate}
\end{corbox}

\begin{proofbox}
\textbf{б)} Понеже $\Omega\cap\varnothing=\varnothing$ (несъвместимо
обединение) и~$\Omega\cup\varnothing=\Omega$, по $\sigma$-адитивност
$\mathbb P(\Omega)=\mathbb P(\Omega)+\mathbb P(\varnothing)
\Rightarrow \mathbb P(\varnothing)=0$.

\textbf{а)} От $\varnothing\subseteq A\subseteq\Omega$ и~монотонността (доказана в~(е) по-долу, независимо от (а)) получаваме
$\mathbb P(\varnothing)\le\mathbb P(A)\le\mathbb P(\Omega)$, т.е.\
$0\le\mathbb P(A)\le 1$.

\textbf{в)} От $\Omega=A\cup\bar A$ (несъвместимо обединение) и~$\sigma$-адитивността: $1=\mathbb{P}(A)+\mathbb{P}(\bar A)
\Rightarrow \mathbb{P}(\bar A)=1-\mathbb{P}(A)$.

\textbf{г)} Директно от $\sigma$-адитивността, приложена към две
несъвместими събития (допълнена с~$\varnothing,\varnothing,\dots$).

\textbf{д)} Пишем $A\cup B = A\cup(B\setminus A)$, като двете части
са несъвместими, следователно $\mathbb{P}(A\cup B)=\mathbb{P}(A)+
\mathbb{P}(B\setminus A)$. От друга страна $B=(B\setminus A)\cup(A\cap
B)$ (несъвместимо разбиване), т.е.\ $\mathbb{P}(B)=\mathbb{P}(B
\setminus A)+\mathbb{P}(A\cap B)$, откъдето $\mathbb{P}(B\setminus A)
=\mathbb{P}(B)-\mathbb{P}(A\cap B)$. Заместваме и~получаваме
$\mathbb{P}(A\cup B)=\mathbb{P}(A)+\mathbb{P}(B)-\mathbb{P}(A\cap B)$.

\textbf{е)} Пишем $B = A\cup(B\setminus A)$ - несъвместимо
обединение, следователно $\mathbb{P}(B)=\mathbb{P}(A)+
\underbrace{\mathbb{P}(B\setminus A)}_{\geq 0}\geq\mathbb{P}(A)$.

\textbf{ж) (непрекъснатост отдолу)} Нека $A_1\subseteq A_2\subseteq
\dots$ и~$A=\bigcup_{n=1}^{\infty}A_n$. Дефинираме $C_1=A_1$,
$C_k=A_k\setminus A_{k-1}$ за $k\geq 2$. Тогава $C_k$ са по двойки
несъвместими и~$A_n=\bigcup_{k=1}^{n}C_k$, $A=\bigcup_{k=1}^{\infty}
C_k$. По $\sigma$-адитивност $\mathbb P(A_n)=\sum_{k=1}^n \mathbb
P(C_k)$ и~$\mathbb P(A)=\sum_{k=1}^\infty \mathbb P(C_k)$. Но сумата
на реда по дефиниция е границата на частичните суми, откъдето
$\lim_{n\to\infty}\mathbb P(A_n)=\mathbb P(A)$.
\end{proofbox}

\begin{proofbox}
\textbf{з) (непрекъснатост отгоре)} Първо разглеждаме частния случай,
в~който редицата се свива до празното множество: $B_1\supseteq
B_2\supseteq\dots\supseteq B_n\supseteq\dots$ и~$\bigcap_{j=1}^{\infty}
B_j=\varnothing$. Трябва да покажем, че
$0=\mathbb P(\varnothing)=\lim_{n\to\infty}\mathbb P(B_n)$.

Всеки елемент на $B_1$ отпада от редицата на някоя стъпка, защото
иначе би бил във всички $B_j$, т.е.\ в~сечението. Затова $B_1$ се
разбива на несъвместими „венци“ $B_j\setminus B_{j+1}$, и~по същата
причина това важи за всяко $B_n$:
\[
  B_1=\bigcup_{j=1}^{\infty}(B_j\setminus B_{j+1}),
  \qquad
  B_n=\bigcup_{j=n}^{\infty}(B_j\setminus B_{j+1}).
\]
По $\sigma$-адитивност
\[
  1\ \ge\ \mathbb P(B_1)=\sum_{j=1}^{\infty}\mathbb P(B_j\setminus B_{j+1})<\infty,
  \qquad
  \mathbb P(B_n)=\sum_{j=n}^{\infty}\mathbb P(B_j\setminus B_{j+1})
  \xrightarrow[n\to\infty]{}0.
\]
Първият ред е сходящ, а $\mathbb P(B_n)$ е неговата остатъчна сума
след $n-1$ члена, която за сходящ ред клони към~$0$. Това е точно
непрекъснатостта в~$\varnothing$.

Общият случай се свежда до частния. Нека $A_1\supseteq A_2\supseteq
\dots$ и~$A=\bigcap_{i=1}^{\infty}A_i$. Полагаме $B_n=A_n\setminus A$.
Тогава $B_n$ намаляват и~$\bigcap_n B_n=\varnothing$, а понеже
$A\subseteq A_n$, имаме $\mathbb P(A_n)=\mathbb P(A)+\mathbb P(B_n)
\to\mathbb P(A)$.

\textbf{и) (субадитивност, чрез „венци“)} Нека
\[
  B_1=A_1,\quad B_2=A_2\setminus A_1,\quad B_3=A_3\setminus(A_2\cup A_1),\quad \dots,\quad B_n=A_n\setminus\bigcup_{i=1}^{n-1}A_i.
\]
Тогава $B_i\cap B_j=\varnothing$ за $i\neq j$ (по конструкция) и~$B_i\subseteq A_i$ за всяко $i$, откъдето по монотонност (е)
$\mathbb{P}(B_i)\leq\mathbb{P}(A_i)$. Освен това
$\bigcup_{i}B_i=\bigcup_i A_i$: включването „$\subseteq$“ е очевидно
($B_i\subseteq A_i$); за „$\supseteq$“ - ако $\omega\in\bigcup_i A_i$,
нека $k=\min\{i\geq 1:\omega\in A_i\}$; тогава $\omega\notin
\bigcup_{i=1}^{k-1}A_i$, т.е.\ $\omega\in B_k\subseteq\bigcup_i B_i$.
Следователно, по $\sigma$-адитивност и~монотонност,
\[
  \mathbb{P}\Big(\bigcup_{i}A_i\Big) = \mathbb{P}\Big(\bigcup_i B_i\Big) = \sum_i \mathbb{P}(B_i) \leq \sum_i \mathbb{P}(A_i).
\]
Крайният случай се получава аналогично (или чрез допълване с~$\varnothing$). \qedbg
\end{proofbox}

\section{Вероятностно пространство}

\begin{defbox}[~(Вероятностно пространство)]
\textbf{Вероятностно пространство} е наредена тройка
$(\Omega,\mathcal{A},\mathbb{P})$, където:
\begin{itemize}[nosep]
  \item $\Omega$ е пространството на елементарните събития - непразно множество от всички възможни резултати от даден случаен експеримент;
  \item $\mathcal{A}$ е $\sigma$-алгебра на събитията върху $\Omega$;
  \item $\mathbb{P}$ е вероятностна мярка, т.е.\ функция $\mathbb{P}:\mathcal{A}\to[0,1]$, удовлетворяваща аксиомите от предната дефиниция.
\end{itemize}
\end{defbox}

\begin{rembox}[Забележка]
Ако първоначалното вероятностно пространство е било
$(\Omega,\mathcal{A},\mathbb{P})$ и~ограничим вниманието си само до
събития, съдържащи се в~дадено фиксирано $B\in\mathcal A$ с~$\mathbb P(B)>0$, естествено е новото пространство от елементарни
събития да стане самото $B$ (вместо цялото $\Omega$), с~$\sigma$-алгебра
$\mathcal A\cap B := \{A\cap B : A\in\mathcal A\}$ (следата на
$\mathcal A$ върху $B$) и~нова вероятностна мярка $\mathbb P(\cdot
\mid B)$ - идея, върху която стъпва условната вероятност в~следващата глава.
\end{rembox}
